\clearpage
\section{Methodology}\label{sec:methods}
\subsection{Preprocessing and motivation}\label{sec:preprocess}
Both branches start with minimal cleanup: replace HTML tags with spaces, decode
HTML entities and normalize whitespace. This prevents formatting artifacts from
becoming features or joining neighboring words. Original text remains available
for inspection. Classical preprocessing then applies the following techniques.

\paragraph{Tokenization.}
spaCy's \code{en_core_web_sm} pipeline divides cleaned text into linguistic
tokens, separating punctuation and contractions. For example, contracted
``didn't'' is split into ``did'' and ``n't''. Tokenization establishes the units
for word-level filtering and feature construction. The linguistic pipeline runs
on the complete sentence before filtering, so its tagger and lemmatizer can use
the original context.

\paragraph{Stop-word removal.}
Tokens flagged as stop words are removed to reduce common function words and
feature volume. However, negation and contrast can determine sentiment, so the
protected set retains \code{not}, \code{no}, \code{never}, \code{neither},
\code{nor}, contracted negation, and contrast markers such as \code{but},
\code{although}, \code{however}, \code{though} and \code{yet}. Punctuation-only
and whitespace tokens are filtered separately. Without the exceptions,
``not good'' could lose its negative cue. Filtering remains a modeling choice
that may also remove useful information.

\paragraph{Lemmatization.}
Contextual lemmas map inflected forms to dictionary forms, for example ``actors''
to ``actor'' and ``enjoyed'' to ``enjoy''. This reduces vocabulary sparsity for
word-based features. Classical tokens are also lowercased; contracted negation
is normalized to \code{not}, and protected words retain normalized surface forms.
The demonstration sentence ``I didn't enjoy this movie, but the actors were
excellent!'' becomes \code{not enjoy movie but actor excellent}.

Batch processing uses \code{nlp.pipe} with batch size 64. Table~\ref{tab:tokens}
summarizes the reduction in token volume. Initial spaCy counts include punctuation;
final counts include retained normalized tokens, so the difference is not due
to stop-word removal alone. No preprocessing ablation was run, and these counts
do not establish an accuracy benefit.

\begin{table}[htbp]\centering
\caption{Mean token counts before and after classical filtering and normalization.}\label{tab:tokens}
\begin{tabular}{lrrr}\toprule
Split & Initial spaCy tokens & Classical tokens & Empty outputs\\\midrule
Training & 272.71 & 111.57 & 0\\
Validation & 268.19 & 109.58 & 0\\
Test & 265.39 & 108.45 & 0\\\bottomrule
\end{tabular}
\end{table}
The transformer uses \code{text_clean} with its own tokenizer, preserving natural
syntax, punctuation and function words before uncased tokenization. Learned
TF-IDF statistics and Word2Vec embeddings use training records only.

\subsection{Classical text representations}
Figure~\ref{fig:mlworkflow} shows the four classical combinations. The same fixed
training, validation and test memberships are used throughout.
\reportfigure{person2_image1.png}{0.98}{Our classical workflow: two representations combined with two linear classifiers.}{fig:mlworkflow}

\subsubsection{TF-IDF}
The implementation uses scikit-learn's \code{TfidfVectorizer} on
\code{text_classical}, with unigram and bigram features,
\code{min_df=2}, \code{max_df=0.95} and \code{sublinear_tf=True}.
For term $t$ in document $d$, the weighting corresponds to
\begin{align}
 \operatorname{tf}(t,d)&=1+\log c(t,d)\quad\text{if }c(t,d)>0,\\
 \operatorname{idf}(t)&=\log\frac{1+N}{1+\operatorname{df}(t)}+1,\\
 \widetilde{x}_{t,d}&=\operatorname{tf}(t,d)\operatorname{idf}(t),\qquad
 x_d=\frac{\widetilde{x}_d}{\|\widetilde{x}_d\|_2}.
\end{align}
Here $N$ is the number of training documents, $c(t,d)$ the term count and
$\operatorname{df}(t)$ the training document frequency. Smoothed IDF and
$L_2$ normalization are the vectorizer defaults. Vocabulary and IDF are fitted
on training only; validation and test use the frozen transformation.

The vectorizer retains its default word analyzer, including the token pattern
that accepts word tokens of at least two characters. Thus its effective tokens
can differ from our spaCy token list, even though no second stop-word
list is applied. Bigrams provide local lexical interactions after preprocessing,
while the sparse representation does not model long-range context.

\subsubsection{Word2Vec and document pooling}
A Gensim Skip-gram Word2Vec model is trained on \code{tokens_classical} from the
training split. Its configuration is 100 dimensions, context window 5, minimum
count 2, \code{sg=1}, 10 epochs, one worker and seed 42. Single-worker training
reduces one source of run-to-run variation. Validation and test do not update
the embedding vocabulary or weights.

For review $d$, let $K_d$ be its token occurrences represented in the learned
vocabulary and $e(w)\in\mathbb{R}^{100}$ a word vector. The document representation is
\begin{equation}
 v_d=\begin{cases}
 |K_d|^{-1}\sum_{w\in K_d} e(w), & |K_d|>0,\\
 \mathbf{0}\in\mathbb{R}^{100}, & |K_d|=0.
 \end{cases}
\end{equation}
Unknown words are omitted from the average. A zero vector is the fallback if
no word is known. Averaging creates a compact fixed-size input but discards
order and may dilute a decisive phrase. Figure~\ref{fig:representations}
contrasts this representation with sparse lexical features.
\reportfigure{person2_image2.png}{0.93}{Conceptual comparison of sparse TF-IDF and mean-pooled Word2Vec representations.}{fig:representations}

\subsection{Classical classifiers}
\paragraph{Logistic Regression (LR).}
Each representation is used to fit a separate linear classifier with
$P(y=1\mid x)=\sigma(w^\top x+b)$ and \code{max_iter=1000}. The regularization
parameter remains at its default $C=1.0$. The classifier returns labels and
positive-class probabilities.

\paragraph{Linear Support Vector Classifier (LinearSVC).}
A separate \code{LinearSVC} is fitted to each representation with default
$C=1.0$. Its linear margin is suitable for high-dimensional sparse inputs.
It outputs labels and decision scores, rather than calibrated probabilities.
For the demo, a sigmoid transforms the margin to $[0,1]$, but this does not
constitute probability calibration. No feature scaling or alternative document
pooling step is added in the recorded implementation.

\subsection{DistilBERT with LoRA}\label{sec:dlmethod}
\reportfigure{person3_image1.png}{0.98}{DistilBERT + LoRA workflow. Encoder weights are frozen; low-rank updates and the classification head are trained. The illustrated percentage uses the recorded base-model reference count.}{fig:dlworkflow}
\subsubsection{Input and long-review handling}
The deep branch reads \code{text_clean}. The tokenizer of
\code{distilbert-base-uncased} lowercases text and splits it into WordPiece units
from a 30,522-entry vocabulary. Subword tokenization reduces dependence on
whole-word vocabulary coverage, but does not guarantee that unknown tokens
never occur. Each input has the form \code{[CLS] review [SEP]}.
Dynamic padding extends sequences to the longest input in each batch, and an
attention mask excludes padding positions.

Maximum lengths of 256 and 512 include the two special tokens, leaving 254 and
510 review wordpieces. Head-only truncation keeps the beginning. The selected
head-plus-tail strategy uses a head fraction of 0.25:
\begin{equation}
 B=L-2,\qquad H=\operatorname{round}(0.25B),\qquad T=B-H.
\end{equation}
At $L=512$, this means the first 128 and last 382 review wordpieces. Shorter
reviews remain intact. The motivation is to retain both opening context and
closing judgments, while acknowledging that evidence in the middle can be lost.
The fraction of training reviews truncated falls from 41.41\% at length 256 to
13.46\% at length 512.

\subsubsection{Encoder, head and low-rank updates}
DistilBERT has six Transformer blocks, hidden size 768 and 12 attention heads.
The hidden state at \code{[CLS]} passes through a 768-to-768
\code{pre_classifier}, ReLU and dropout, followed by a 768-to-2 \code{classifier}.
Both classification layers are randomly initialized and trained in full. The
code explicitly preserves both through PEFT's \code{modules_to_save}.

For an adapted projection, LoRA learns
\begin{equation}
 h=W_0x+\frac{\alpha}{r}BAx,
 \qquad A\in\mathbb{R}^{r\times d_{\mathrm{in}}},\quad
 B\in\mathbb{R}^{d_{\mathrm{out}}\times r}.
\end{equation}
The pretrained $W_0$ remains frozen. The selected configuration adapts the query
and value projections, \code{q_lin} and \code{v_lin}, in all six blocks with
$r=8$, $\alpha=16$ and LoRA dropout 0.1. Its trainable components are listed in
Table~\ref{tab:params} and highlighted in Figure~\ref{fig:dlworkflow}.

\begin{table}[htbp]\centering
\caption{Trainable parameters of the selected deep pipeline.}\label{tab:params}
\begin{tabularx}{\linewidth}{@{}Ylr@{}}\toprule
Component & Calculation & Parameters\\\midrule
Query/value LoRA adapters & $6\times2\times8\times(768+768)$ & 147,456\\
\code{pre_classifier} & $768\times768+768$ & 590,592\\
\code{classifier} & $768\times2+2$ & 1,538\\\midrule
Total trainable & & 739,586\\\bottomrule
\end{tabularx}
\end{table}
The script reports a base sequence-classifier count of 66,955,010, measured
before PEFT wrapping. The trainable count is 1.1046\% of that reference count;
this percentage uses the script's base-model denominator rather than claiming
to count every additional adapter and wrapper copy in the adapted object.

\FloatBarrier
